% !TeX root = ../tfg.tex
% !TeX encoding = utf8
%
%*******************************************************
% Capítulo 5. Métodos de detección post hoc
%*******************************************************

\chapter{Métodos de detección post hoc}\label{cap:metodos}

% Notación: logits f(x) \in \R^C, características h(x) \in \R^D (D = 768), pesos de la última capa
% W \in \R^{C \times D} y sesgo b, de modo que f(x) = W h(x) + b.

\section{Métodos basados en logits}\label{sec:metodos-logits}

\pendiente{métodos basados en logits}

\section{Métodos basados en la modificación de activaciones}\label{sec:metodos-activaciones}

\pendiente{métodos basados en la modificación de activaciones}

\section{Métodos basados en características y distancias}\label{sec:metodos-caracteristicas}

\pendiente{métodos basados en características}

\section{STOOD-X}\label{sec:metodo-stoodx}

\pendiente{STOOD-X}

\endinput
